% ================================================================================================ capítulo 13
\chapter{Por dentro: el código}\label{cap:codigo}

\begin{resumen}
\textbf{En este capítulo:} cómo está organizado el código de TurbuLENS y los fragmentos que hacen las mediciones,
copiados del código fuente de esta versión con sus números de línea reales, cada uno con su explicación. Sirve para
revisar el método, citarlo con precisión o mejorarlo.
\end{resumen}

TurbuLENS está escrito en Python 3. La interfaz usa Qt (PySide6); los cálculos, NumPy; los videos se leen con FFmpeg a
través de PyAV, y los archivos del satélite, con pyfive, un lector de HDF5 escrito en Python puro que viene incluido. El
código es abierto (licencia PolyForm Noncommercial, apéndice~\ref{sec:licencia}) y está en
\url{https://github.com/DiegoCagua/TurbuLENS}.

\section{Cómo está organizado}

La regla principal es que \textbf{el motor no sabe que existe una ventana}: todo lo que mide, calcula o escribe archivos
está en \ruta{turbulens/core/} y no importa Qt. Así se puede probar sin interfaz, usar desde un script y revisar sin
distraerse con botones.

\begin{center}\small
\begin{tabularx}{\linewidth}{@{}>{\ttfamily\color{acentoosc}}l Y@{}}
\toprule
\normalfont Archivo & Qué hace\\
\midrule
core/video.py & Lee videos con PyAV: información, cuadros en bloques, búsqueda por cuadro clave.\\
core/pipeline.py & El procesamiento: recorte, calibración, seguimiento y detección de cada cuadro.\\
core/detect.py & La detección de los dos puntos (k-medias con dos grupos y centroides).\\
core/dimm.py & Del movimiento diferencial al seeing (Tokovinin).\\
core/stats.py & Estadística: resumen, recorte robusto, histogramas, correlación.\\
core/devices.py & Telescopios, cámaras y máscaras del OAN y propios; escala de placa.\\
core/calframes.py & Calibración guiada por las cabeceras: clasificación, plan, maestros.\\
core/calibration.py & Combinación de imágenes y corrección con dark y flat.\\
core/starfit.py & Ajuste de una gaussiana 2D a una estrella (Levenberg–Marquardt).\\
core/goes.py & Satélite GOES: búsqueda de archivos, lectura parcial, proyección, lectura de un punto.\\
core/satseeing.py & Seeing desde el satélite: $F$, tendencia, Sol, ángulo del satélite.\\
core/mapdata.py & Costas, fronteras y departamentos de Natural Earth, proyectados sobre el recorte.\\
core/fits.py, imageio.py & Lectura y escritura de FITS, PNG de 16 bits, TIFF y NumPy.\\
core/export.py, writers.py & Exportación de datos y escritura de cuadros.\\
core/history.py & El historial en SQLite.\\
charts.py, reports.py & Gráficas dibujadas con Qt (pantalla, PNG, SVG, PDF) e informes.\\
gui/ & La interfaz: una página por pantalla, el visor, las animaciones, los temas y los idiomas.\\
locales/ & Traducciones al inglés y al portugués (la clave es el texto en español).\\
\bottomrule
\end{tabularx}
\end{center}

\begin{figure}[H]
\centering
\begin{tikzpicture}[font=\sffamily\scriptsize,>=Stealth,
  caja/.style={draw=acento,fill=suave,rounded corners=3pt,align=center,minimum height=1cm,text width=2.45cm},
  dato/.style={draw=naranja,fill=crema,rounded corners=3pt,align=center,minimum height=0.8cm,text width=2.45cm}]
  \node[dato] (v) at (0,0) {Video\\\ttfamily video.py};
  \node[caja] (p) at (3.2,0) {Recorte, calibración\\y seguimiento\\\ttfamily pipeline.py};
  \node[caja] (d) at (6.4,0) {Dos puntos\\por cuadro\\\ttfamily detect.py};
  \node[dato] (m) at (9.6,0) {Mediciones\\(arreglo por cuadro)};
  \node[caja] (a) at (9.6,-1.9) {Seeing, $r_0$,\\controles\\\ttfamily dimm.py};
  \node[dato] (r) at (6.4,-1.9) {Resultados};
  \node[caja,fill=white] (e) at (3.2,-1.9) {Exportar, informe\\\ttfamily export.py};
  \node[caja,fill=white] (h) at (0,-1.9) {Historial\\\ttfamily history.py};
  \draw[->,gris] (v) -- (p); \draw[->,gris] (p) -- (d); \draw[->,gris] (d) -- (m); \draw[->,gris] (m) -- (a);
  \draw[->,gris] (a) -- (r); \draw[->,gris] (r) -- (e); \draw[->,gris] (e) -- (h);
  \draw[->,gris,dashed] (r.south) to[bend left=18] (h.south);
\end{tikzpicture}
\caption{El camino de un video hasta el seeing.}
\end{figure}

Los fragmentos que siguen se copian del código de esta versión con un script (\ruta{scripts/manual\_codigo.py}), así que
los números de línea son los del archivo. Los comentarios del código están en español, como el resto del proyecto.

% ------------------------------------------------------------------------------------------------ detección
\section{Detección de los dos puntos}\label{sec:codigo_deteccion}

La función \codigo{\_detect} recibe un bloque de cuadros con forma \codigo{(n, alto, ancho)} y devuelve, para cada
cuadro, las posiciones de los dos puntos y si el cuadro es válido. Nada se recorre con un bucle de Python: cada línea
opera sobre todos los cuadros a la vez.

\fragmento{umbral}{Fondo, ruido y umbral}

Para cada cuadro, el fondo \codigo{bg} es la mediana y el ruido, la MAD por 1,4826 (con un mínimo de 0,5 cuentas, para
cámaras de 8 bits con fondo perfectamente plano). \codigo{s3} es la imagen suavizada con una caja de 3~×~3 hecha con sumas
acumuladas. El umbral \codigo{thr} es el mayor entre el 15\,\% del pico suavizado y $\frac43$ del ruido; los píxeles de
luz (\codigo{fg}) deben superarlo en la imagen suavizada y en la original, y su peso \codigo{wgt} es cuánto lo superan.
Restar el umbral hace que los bordes, más ruidosos, pesen poco en el centroide.

\fragmento{centroides}{k-medias con dos grupos}

\codigo{centroids} asigna cada píxel al centro más cercano (\codigo{A = d1 <= d2}) y calcula el centroide ponderado de
cada grupo: $\bar x = \sum w_i x_i / \sum w_i$. Si un grupo se queda sin píxeles conserva su centro anterior. El bucle
repite hasta que ningún centro se mueva más de una milésima de píxel o hasta \codigo{cfg.iterations} (8) veces. Los
argumentos \codigo{lim1} y \codigo{lim2} se usan en una segunda pasada, que deja fuera los píxeles a más de $3r + 3$
píxeles de cada centro.

\fragmento{valido}{Cuándo un cuadro es válido}

Cada condición es un arreglo booleano de un elemento por cuadro, y el cuadro es válido si se cumplen todas: los dos
grupos tienen luz; sus picos superan seis veces el ruido (\codigo{snr\_min}); el flujo del más débil es al menos el
12\,\% del otro (\codigo{second\_min\_ratio}); están separados más de 2,2 veces el mayor radio de giro; en el punto
medio la imagen suavizada cae por debajo del 60\,\% del pico menor (\codigo{dip}: hay un valle entre los dos), y ninguno
toca el borde (\codigo{inside}).

\subsection*{El seguimiento}

Dentro del bucle que recorre los cuadros del video, cuando \campo{Seguir a la estrella} está activo:

\fragmento{seguimiento}{La región sigue a la estrella}

\codigo{light\_center} da el centro de luz del recorte. La región se mueve solo la mitad del camino hacia él: es un filtro
que sigue la deriva lenta del telescopio sin copiar el temblor rápido de la turbulencia. Si en cinco cuadros seguidos no
hay estrella, \codigo{locate\_pair} la busca en el cuadro completo.

% ------------------------------------------------------------------------------------------------ DIMM
\section{Del vector al seeing}

\fragmento{ejebase}{Orientación, atípicos y componentes}

Los dos puntos pueden intercambiarse entre cuadros. Para orientar todos los vectores igual, se calcula la matriz de
segundos momentos de las separaciones y su vector propio principal (\codigo{np.linalg.eigh}): es la dirección de la línea
de base. Los vectores que apuntan al revés se invierten (\codigo{sgn}). Después se descartan los atípicos con
\codigo{sigma\_clip\_mask} y se proyecta cada vector sobre la dirección media (\codigo{ux, uy}): \codigo{lon} es la
componente longitudinal y \codigo{tra}, la transversal (el producto con el vector perpendicular $(-u_y, u_x)$).

\fragmento{tokovinin}{Coeficientes de Tokovinin, \texorpdfstring{$r_0$}{r0} y seeing}

Son, línea por línea, las ecuaciones \eqref{eq:kl}, \eqref{eq:kt}, \eqref{eq:r0dimm} y \eqref{eq:seeing}. Las funciones
devuelven \codigo{nan} (\emph{not a number}) si les llega un dato imposible, en lugar de un número sin sentido.

\fragmento{seeingdimm}{El seeing de un análisis}

\codigo{sl} y \codigo{st} son las desviaciones estándar en píxeles. \codigo{conv} pasa píxeles al cuadrado a radianes al
cuadrado con la escala de placa. El seeing estimado es el promedio de los dos, lo mismo que $r_0$; con la altura de la
estrella entre 5° y 90° se calcula el seeing en el cenit.

\fragmento{sigmaclip}{Recorte robusto de atípicos}

Hasta \codigo{iters} veces: mediana, MAD escalada y fuera lo que esté a más de \codigo{k} de ellas. Se detiene si ya no
sale ningún dato o si quedan menos de tres.

\section{Escala de placa}

\fragmento{escala}{Píxel efectivo de una cámara}
\fragmento{escalaplaca}{Escala de placa}

Si la cámara reduce el sensor al grabar (\codigo{video}) y el análisis es de un video, el píxel es el ancho del sensor
entre el ancho del video. La constante 206,265 sale de pasar radianes a segundos de arco (206\,265) con el píxel en
micrómetros y la focal en milímetros ($10^{-3}$).

% ------------------------------------------------------------------------------------------------ calibración
\section{Calibración}

\fragmento{combinar}{Combinar una serie en un maestro}

\codigo{stack} tiene forma \codigo{(n, alto, ancho)} y todo se calcula a lo largo del primer eje, píxel por píxel. En la
media con recorte, \codigo{ok} marca los valores a menos de \codigo{sigma} desviaciones robustas de la mediana de su
píxel; se suman solo esos y se dividen entre cuántos son. Para series que no caben en la memoria, \codigo{calframes}
llama a esta función por franjas de filas.

\fragmento{darkescalado}{El dark que se resta}

El plan decide para cada grupo un modo: el dark de la misma exposición (\codigo{"igual"}), solo el bias, nada, o el
\codigo{"escalado"}, que aplica la ecuación de la sección~\ref{sec:dark_escalado} con el factor $t_\text{luz}/t_\text{dark}$
guardado en \codigo{c.scale}. Los maestros se calculan una vez y se guardan en una caché.

\fragmento{corregir}{Corregir una imagen}

Resta el dark (o, si no hay, el bias) y divide por el flat. Si un maestro no mide lo mismo que la imagen, el mensaje de
error dice los dos tamaños.

% ------------------------------------------------------------------------------------------------ estrella
\section{Ajuste de una estrella}\label{sec:codigo_gauss}

\fragmento{gauss}{Levenberg–Marquardt}

\codigo{p} son los seis parámetros $(A, x_0, y_0, \sigma_x, \sigma_y, B)$, que arrancan en los momentos de la imagen.
\codigo{\_model\_jac} da el modelo y su matriz jacobiana (las derivadas respecto a cada parámetro). En cada iteración se
resuelve $(J^TJ + \lambda\,\operatorname{diag}J^TJ)\,\Delta p = J^T r$: con $\lambda$ pequeño es un paso de Gauss–Newton,
con $\lambda$ grande un paso corto por el gradiente. Si el paso mejora el ajuste, se acepta y $\lambda$ se divide entre 3;
si no, se multiplica por 4 y se intenta otro. Los anchos se mantienen positivos y de al menos 0,2~px. Termina cuando ningún
parámetro cambia más de una parte en diez millones.

% ------------------------------------------------------------------------------------------------ satélite
\section{Satélite}

\fragmento{satformulas}{Función de estabilidad y seeing}

\codigo{big\_l} es $\Lambda(\theta)$; \codigo{cn2}, la integral de $C_n^2$; \codigo{fwhm}, la fórmula del seeing en
radianes por \codigo{ARCSEC} para tenerlo en segundos de arco, y \codigo{r0}, la ecuación~\eqref{eq:r0} con $k = 2\pi/\lambda$.

\fragmento{tendencia}{La tendencia de la noche}

\codigo{np.polyfit} con grado 1 es la recta de mínimos cuadrados. \codigo{use} marca las horas que cumplen el criterio del
Sol; la tendencia se evalúa en todas las horas, también en las excluidas, para dibujarla completa.

\fragmento{cenitsat}{Ángulo cenital del satélite}

\fragmento{sol}{Altura del Sol}

Las fórmulas del \emph{Almanaque Astronómico}: \codigo{n} son los días desde el mediodía del 1 de enero de 2000 (J2000);
\codigo{L} y \codigo{g}, la longitud media y la anomalía media del Sol; \codigo{lam}, su longitud eclíptica; \codigo{eps},
la oblicuidad de la eclíptica; \codigo{ra} y \codigo{dec}, sus coordenadas ecuatoriales; \codigo{gmst}, el tiempo sidéreo
de Greenwich, y \codigo{ha}, el ángulo horario en el sitio. La última línea es la altura sobre el horizonte.

\fragmento{proyeccion}{De latitud y longitud a ángulos de barrido}

La transcripción de las fórmulas de la guía de usuario de GOES-R. Una prueba compara el resultado con el ejemplo numérico
de esa guía.

\subsection*{Lectura parcial}\label{sec:codigo_rango}

\fragmento{rango}{Un bloque de un archivo remoto}

El archivo remoto se divide en bloques de 128~KB. Para leer el bloque \codigo{i} se pide al servidor el rango de bytes
correspondiente; la respuesta debe ser \codigo{206 Partial Content}, y si no lo es se lanza \codigo{RangeNotSupported}
para que TurbuLENS descargue el archivo completo. Los bloques se guardan en una caché de hasta 96 (12~MB): leer la
cabecera del HDF5 y los bloques comprimidos de un recorte vuelve muchas veces a los mismos.

\fragmento{carpeta}{El nombre de la carpeta de cada evaluación}

% ------------------------------------------------------------------------------------------------ pruebas
\section{Las pruebas}

TurbuLENS tiene casi cincuenta pruebas automáticas en \ruta{tests/}, que corren en GitHub Actions con cada cambio; además, los ejecutables de Windows, macOS y Linux se prueban antes de publicarse.
La mayoría usan datos sintéticos de valores conocidos, de modo que se sabe cuál debe ser la respuesta:

\begin{table}[H]
\centering\small
\begin{tabularx}{\linewidth}{@{}>{\ttfamily}l Y@{}}
\toprule
\normalfont Archivo & Qué comprueba (algunos ejemplos)\\
\midrule
test\_core.py & Que la detección encuentra los centroides con precisión de centésimas de píxel, con puntos pequeños,
grandes y en anillo; que rechaza los cuadros con un solo punto; los coeficientes de Tokovinin y un valor de referencia del
seeing; que el análisis recupera una $\sigma$ conocida y cancela un movimiento común; el recorte de atípicos; el ajuste de
una estrella; que sin detecciones las gráficas quedan vacías.\\
test\_calframes.py & La clasificación por cabecera, nombre y carpeta; que el plan empareja exposiciones y filtros y escala
darks; que la calibración quita un píxel caliente y el viñeteo; las series grandes por franjas; el aviso de flats
saturados; la escala de placa de los equipos del OAN.\\
test\_modes.py & La proyección contra el ejemplo de la guía de GOES-R; la elección del satélite por fecha; la lectura de
archivos HDF5 comprimidos; que la lectura parcial por HTTP da los mismos valores que el archivo local; las fórmulas del seeing
satelital y el nombre de las carpetas.\\
test\_gui.py, test\_i18n.py & La interfaz completa sin pantalla: entrar a cada modo, procesar un video, las lecciones con
\emph{Mi video}; que los tres idiomas tienen los mismos textos.\\
\bottomrule
\end{tabularx}
\end{table}

\section{Correr TurbuLENS desde el código}

\begin{tcolorbox}[enhanced,colback=codigofondo,colframe=grisclaro,boxrule=0.6pt,arc=1.5mm,fontupper=\ttfamily\small,
  left=4mm,top=2mm,bottom=2mm]
git clone https://github.com/DiegoCagua/TurbuLENS\\
cd TurbuLENS\\
pip install -r requirements.txt\\
python -m turbulens\hfill\normalfont\color{gris}abre la aplicación\\
\ttfamily python -m turbulens --autoprueba\hfill\normalfont\color{gris}prueba sin ventana\\
\ttfamily python -m pytest\hfill\normalfont\color{gris}las pruebas\\
\ttfamily pyinstaller packaging/turbulens.spec\hfill\normalfont\color{gris}el ejecutable de tu sistema
\end{tcolorbox}

Los textos de la interfaz están en español en el código y pasan por la función \codigo{tr()}, que busca su traducción en
\ruta{locales/en.json} y \ruta{locales/pt.json}. Para agregar un idioma basta con un archivo más con las mismas claves.
Las capturas de este manual se generan con \ruta{scripts/capturas\_manual.py} y los fragmentos de código con
\ruta{scripts/manual\_codigo.py}; después de cambiar la aplicación, se corren los dos y se compila el manual otra vez.
